\section{Débruiter sans effacer les structures}
\label{sec:context}

Une image acquise contient des variations d'intensité qui proviennent
à la fois de la scène et du processus d'observation. Certaines sont
utiles : une frontière entre deux objets, une transition d'éclairage
ou un détail fin. D'autres constituent du bruit. Débruiter consiste à
réduire ces dernières sans effacer les premières. Ce compromis est
difficile parce qu'un calcul local voit seulement les valeurs et leurs
variations ; il ne connaît pas leur origine.

Une moyenne locale atténue efficacement des fluctuations indépendantes,
mais mélange également les intensités de part et d'autre d'une frontière.
Une évolution par EDP organise ce lissage en une famille d'états : le
temps de diffusion règle l'échelle de transformation. L'image initiale
est l'observation, et chaque instant définit une reconstruction possible.
Il ne s'agit pas de mesurer une durée physique d'acquisition. L'EDP
constitue ici un modèle de traitement, dont la pertinence dépend de
l'image, du bruit et du choix d'arrêt.

La chaleur offre une référence simple, linéaire et analysable. Elle
échange de l'intensité entre voisins sans adapter le coefficient aux
variations. Perona--Malik propose au contraire un frein dépendant des
pentes locales \cite{perona1990}. Cette idée motive une comparaison,
mais pas une conclusion préalable : une grande différence peut aussi
être produite par le bruit, et la conductance peut alors préserver ce
qu'il faudrait éliminer.

\subsection{Question et nature du travail}
La question étudiée est :
\begin{quote}
Dans quelles conditions la diffusion dépendant du gradient modifie-t-elle
le compromis entre réduction du bruit et préservation des contours,
par rapport à la diffusion linéaire ?
\end{quote}
Le mot « conditions » inclut le niveau de bruit, le paramètre $K$, la
conductance, la durée d'évolution et la nature de l'image. La réponse
sera limitée à la grille expérimentale observée. Les figures apportent
une lecture des transitions, tandis que les métriques quantifient des
écarts à une référence connue. Aucune de ces lectures ne remplace toutes
les autres.

Le projet est un travail de cours de niveau M2, non une recherche
nouvelle. Son apport pédagogique est de relier la modélisation aux
équations, les équations à un schéma exact et les sorties à une
interprétation traçable. Les hypothèses et limites sont donc aussi
importantes que les scores. La stabilité d'un calcul, sa fidélité à un
modèle et l'efficacité de son débruitage sont trois questions différentes.
Le périmètre comprend la chaleur et les deux lois de conductance dans
une variante directionnelle conservative. Il exclut le défloutage
inverse, la segmentation, les modèles variationnels supplémentaires et
les méthodes apprises. Les perspectives sont décrites comme telles,
sans leur attribuer de simulations réalisées.

\subsection{Une comparaison et non un concours présupposé}
L'image propre est disponible parce que les données sont synthétiques.
Elle permet de mesurer les reconstructions et d'illustrer un compromis
qui serait plus difficile à quantifier sur une acquisition inconnue.
Elle n'est utilisée par aucun solveur. En revanche, retenir après coup
le meilleur réglage selon une métrique utilise cette information ; le
rapport nomme ce choix « oracle exploratoire ».

Le bruit non traité reste une méthode de référence à part entière.
Une évolution peut améliorer une mesure et en détériorer une autre,
ou devenir moins bonne que l'observation avec un arrêt trop tardif.
Le protocole conserve donc tous les réglages prédéfinis et les cas
défavorables. Cette organisation évite de confondre une illustration
réussie avec une conclusion générale.
